\documentclass[../main.tex]{subfiles}

\begin{document}
    \section{Tensors}
    \begin{definition}[Tensors]\label{def:tensor-finite}
        If \(\dim \manifold[M] \neq \infty\), the tensors of type \((r, s)\) 
        \(T_p^{r, s}\manifold[M]\) are all the linear functions 
        \[
        f : \bigtimes^r T_p^*\manifold[M] \times \bigtimes^s T_p\manifold[M] 
        \to \mathbb{R}.
        \]
        I.e., it eats \(r\) covectors and \(s\) vectors.
    \end{definition}
    \begin{theorem}[Dimensions of General Tensor Space]\label{thm:dim-tensor}
        The dimension of \(T_p^{r, s}\manifold[M]\) is \(m^rm^s\). In particular, 
        a basis for the space is, 
        \[
        \bigotimes_{1 \leq \mu_1 \cdots \mu_r \leq m} \tvec{\mu_i}{p} 
        \otimes \bigotimes_{1 \leq \nu_1 \cdots \nu_s \leq m} (dx^{\nu_i})_{p} 
        \]
    \end{theorem}
    \begin{remark}
        For a detailed proof, see Hoffman.
    \end{remark}
    \begin{theorem}[Transformation Properties of Tensor Fields]\label{thm:tensor-transformation}
        Given a manifold \(\manifold[M]\) of dimension \(m\), choose two overlapping 
        charts \(\phi : U \to V\) with local coordinates \(x^1, \dots, x^m\), 
        \(\phi' : U' \to V'\) with \({x'}^1, \dots, {x'}^m\).
        A tensor field \(T\) of type \((r, s)\) with local representation on \(U\) given by 
        \[
        T\indices{^{\mu_1}^{\dots}^{\mu_r}_{\nu_1}_{\dots}_{\nu_s}}
        \tfld{\mu_1}\otimes\dots\otimes\tfld{\mu_r}\otimes
        dx^{\nu_1}\otimes\dots\otimes dx^{\nu_s}
        \]
        transforms like 
        \[
        T\indices{^{\mu_1'}^{\dots}^{\mu_r'}_{\nu_1'}_{\dots}_{\nu_s'}}
        =T\indices{^{\mu_1}^{\dots}^{\mu_r}_{\nu_1}_{\dots}_{\nu_s}}
        \frac{\partial {x'}^{\mu_1'}}{\partial x^{\mu_1}}\dots
        \frac{\partial {x'}^{\mu_r'}}{\partial x^{\mu_r}}
        \frac{\partial x^{\nu_1}}{\partial {x'}^{\nu_1'}}\dots
        \frac{\partial x^{\nu_1}}{\partial {x'}^{\nu_s'}}
        \]
    \end{theorem}
    \begin{definition}[Pullback of Covariant Tensor]\label{def:covariant-tensor-pullback}
        The pullback of a general tensor field \(\tau\) by \(f : \manifold[M] \to 
        \manifold[N]\) is defined as 
        \[
        (f^*\tau)(X_1, \dots, X_k) := (\tau \circ f)(f_*X_1, \dots, f_*X_k).
        \]
    \end{definition}
    \subsection{The Metric Tensor}
    \subsubsection{Riemannian Metric}
    \begin{definition}[Riemannian Metric]\label{def:riemannian-metric}
        A Riemannian metric on a manifold \(\manifold[M]\) is a totally symmetric, 
        smooth \((2, 0)\)-tensor field on \(\manifold[M]\) that is positive definite 
        at every point \(p \in \manifold[M]\). Locally, 
        \[
        g = g\indices{_i_j}dx^i\otimes dx^j.
        \]
    \end{definition}
    \begin{definition}[The Euclidean Metric]\label{def:euclidean-metric}
        On the Euclidean space \(\mathbb{R}^m\), there is a natural metric called the 
        Euclidean metric, which is locally defined as 
        \[
        \eucMetric=\delta\indices{_i_j}dx^i\otimes dx^j.
        \]
    \end{definition}
    \begin{definition}[The Product Metric]\label{def:product-metric}
        Given two Riemannian manifolds \((\manifold[M], g_{\manifold[M]}), 
        (\manifold[N], g_{\manifold[N]})\), we can form a metric \(\hat{g}\) on the 
        product manifold \(\manifold[M]\times\manifold[N]\) by 
        \[
        \hat{g}((v_1, v_2), (w_1, w_2)) = g_{\manifold[M]}(v_1, w_1) + 
        g_{\manifold[N]}(v_2, w_2).
        \]
    \end{definition}
    \subsubsection{Existence of Riemannian Metric}
    \begin{theorem}[Existence of Riemannian Metric]\label{thm:riemannian-metric-existence}
        Every smooth manifold admits a Riemannian metric \cref{def:riemannian-metric}.
    \end{theorem}
    \begin{proof}
        Let the manifold have coordinate charts \(\phi_j, j \in J\). Let a smooth 
        partition of unity \(\rho_j, j \in J\) subordinate to \(U_j\). Then we can 
        form a smooth tensor field 
        \[
        g = \sum_{j}^{} \rho_j\phi_j^*g_j,
        \]
        where \(g_j\) is the Euclidean metric on \(V_j\). It is obviously symmetric. 
        Since the partition of unity is positive and every \(g_j\) is positive 
        definite, \(g\) is positive definite also.
    \end{proof}
    \newpage
    \subsubsection{Pullback Metric}
    \begin{definition}[The Pullback Metric]\label{def:pullback-metric}
        Let \(f : \manifold[M] \to \manifold[N]\) be \(C^\infty\), where 
        \(\manifold[N]\) is a Riemannian manifold equipped with Riemannian metric 
        \(g\). Then \(f^*g\) is a Riemannian metric on \(\manifold[M]\) iff \(f\) 
        is a smooth immersion.
    \end{definition}
    \begin{theorem}[A Useful Mnemonic for the Pullback Metric]\label{thm:mnemonic-pullback-metric}
        Let \(f : \manifold[M] \to \manifold[N]\) be a smooth immersion so that 
        \cref{def:pullback-metric} can be defined. Then 
        \[
        (f^*g)=(g\indices{_i_j}\circ f)d(y_i\circ f)^2\otimes d(y_j\circ f)^2,
        \]
        where \(d(y_i \circ f)\) expands like \(\frac{\partial (y_i\circ f)}
        {\partial x^j}dx^j\), just like the exterior derivative.
    \end{theorem}
    \begin{proof}
        \[
        \begin{aligned}
            f^*g(X_1, X_2)&=(g\circ f)(f_*X_1, f_*X_2) \\
            &=(g\indices{_i_j}\circ f)dy^i(f_*X_1)dy^j(f_*X_2) \\
            &=(g\indices{_i_j}\circ f)dy^i(\tfld{\nu_1}\frac{\partial f^{\nu_1}}
            {\partial x^{\mu_1}}X_1^{\mu_1})dy^j(\tfld{\nu_2}\frac{\partial f^{\nu_2}}
            {\partial x^{\mu_2}}X_2^{\mu_2}) \\
            &=(g\indices{_i_j}\circ f)\frac{\partial f^i}{\partial x^{\mu_1}}
            \frac{\partial f^j}{\partial x^{\mu_2}}X_1^{\mu_1}X_2^{\mu_2}.
        \end{aligned}
        \]
        So 
        \[
        (f^*g)\indices{_\mu_\nu}=(g\indices{_i_j}\circ f)\frac{\partial f^i}
        {\partial x^{\mu_1}}\frac{\partial f^j}{\partial x^{\mu_2}}.
        \]
        But 
        \[
        d(y^i\circ f) := \frac{\partial (y^i\circ f)}{\partial x^j}dx^j.
        \]
        Tensor product all those completes the proof.
    \end{proof}
    \subsubsection{Musical Isomorphism}
    \begin{definition}[Musical Isomorphism]\label{def:musical-isomorphism}
        \[
        \begin{aligned}
            (X^\flat)\indices{_\mu}&=g\indices{_\mu_\nu}X^\nu \\
            (\omega^\sharp)\indices{^\mu}&=g\indices{^\mu^\nu}\omega_\nu,
        \end{aligned}
        \]
        where \(g\indices{^\mu^\nu}:=g ^{-1}_{\mu\nu}\).
    \end{definition}
\end{document}